\documentclass[11pt]{article}
\usepackage[margin=1in]{geometry}
\usepackage{amsmath,amssymb}
\usepackage{enumitem}
\usepackage[T1]{fontenc}

\title{COMP2300 Set 2 Practice Exam -- Answer Key}
\author{}
\date{}

\begin{document}
\maketitle

\section*{Part A: Multiple Choice}
\begin{enumerate}[label=\textbf{A\arabic*.}]
  \item \textbf{C}
  \item \textbf{B}
  \item \textbf{B}
  \item \textbf{B}
  \item \textbf{B}
  \item \textbf{C}
  \item \textbf{B}
  \item \textbf{B}
  \item \textbf{B}
  \item \textbf{C}
\end{enumerate}

\section*{Part B: Computational and Design Solutions}
\begin{enumerate}[label=\textbf{B\arabic*.}]
  \item \textbf{Truth table from specification}\\
  Output is 1 for inputs with at least two 1s, except 101 is forced to 0.
  \[
  \begin{array}{c c c|c}
  A & B & C & Y\\\hline
  0&0&0&0\\
  0&0&1&0\\
  0&1&0&0\\
  0&1&1&1\\
  1&0&0&0\\
  1&0&1&0\\
  1&1&0&1\\
  1&1&1&1
  \end{array}
  \]
  Canonical SOP: \(Y = A'BC + ABC' + ABC = \Sigma m(3,6,7)\).\\
  Simplify: \(ABC'+ABC=AB\), so \(Y=AB + A'BC = B(A + A'C)=B(A+C)=AB+BC\).

  \item \textbf{SOP to shorthand and implementation}\\
  Given \(F=A'B'C + A'BC + AB'C + ABC'\).\\
  Minterms are \(m_1, m_3, m_5, m_6\), so \(F=\Sigma m(1,3,5,6)\).\\
  This is canonical SOP (sum of minterms), but not guaranteed minimal.\\
  Two-level implementation: four 3-input AND gates feeding one 4-input OR gate (with needed input inversions).

  \item \textbf{2:1 MUX derivation and evaluation}\\
  (a) Behavior gives \(Y=S'D_0 + SD_1\).\\
  (b) Substitute \(D_0=A, D_1=B', S=C\): \(Y=C'A + CB'\).\\
  (c) For \((1,1,0)\): \(Y=1\cdot1 + 0\cdot0=1\).\\
  For \((0,1,1)\): \(Y=0\cdot0 + 1\cdot0=0\).

  \item \textbf{Full adder and 4-bit ripple adder}\\
  (a) Full adder equations:
  \[
  S=A\oplus B\oplus C_{in},\qquad
  C_{out}=AB + C_{in}(A\oplus B).
  \]
  (b) \(1011_2 + 0110_2 = 10001_2\).\\
  4-bit sum output = \(0001_2\), carry-out = 1.

  \item \textbf{Delay analysis}\\
  (a) Ripple-carry delay: \(t_{RCA}(N)=N\cdot t_{FA}\).\\
  (b) For \(N=16\): \(t=16\times14=224\) ps.\\
  (c) CLA computes carry signals with parallel lookahead logic (generate/propagate), reducing linear carry-chain dependence.

  \item \textbf{2:4 decoder equations}\\
  \[
  D_0=A'B',\quad D_1=A'B,\quad D_2=AB',\quad D_3=AB.
  \]
  For \(AB=10\): \((D_0,D_1,D_2,D_3)=(0,0,1,0)\).

  \item \textbf{Priority circuit}\\
  Highest active request wins:
  \[
  OUT0=IRQ0,
  \]
  \[
  OUT1=IRQ0'\,IRQ1,
  \]
  \[
  OUT2=IRQ0'\,IRQ1'\,IRQ2.
  \]
  Cases:\\
  i. \((1,1,1)\to(1,0,0)\)\\
  ii. \((0,1,1)\to(0,1,0)\)\\
  iii. \((0,0,1)\to(0,0,1)\)\\
  iv. \((0,0,0)\to(0,0,0)\)

  \item \textbf{ALU add/sub and flags}\\
  4-bit two's complement range is \([-8,7]\).\\
  (a) Add: \(0111(+7)+0011(+3)=1010\). Mathematical result is +10 (out of range), so overflow.\\
  Result bits \(=1010\). Flags: \(N=1, Z=0, C=0, V=1\).\\
  (b) Subtract: \(1000(-8)-0001(+1)=-9\) (out of range).\\
  In hardware: \(1000 + (1110 + 1)=1000+1111=1\,0111\).\\
  Result bits \(=0111\). Flags: \(N=0, Z=0, C=1, V=1\).

  \item \textbf{Boolean simplification and De Morgan}\\
  (a) \(Y=AB+AB'=A(B+B')=A\).\\
  (b) \(X=A'BC + ABC + AB'C = C(A'B+AB+AB')\).\\
  Since \(AB+AB'=A\), we get \(X=C(A'B + A)=C(A + B)\).\\
  (c) \((A+B+C)' = A'B'C'\).

  \item \textbf{Tri-state buffer reasoning}\\
  (a) To avoid contention, at most one enable can be 1 at a time (never \(E_{CPU}=E_{MEM}=1\)).\\
  (b) If both enables are 0, bus is high-impedance (floating): \(Z\).
\end{enumerate}

\section*{Part C: Past Exam Questions -- Reference Solutions}
\begin{enumerate}[label=\textbf{C\arabic*.}]
  \item \textbf{2023\_Final\_Exam.pdf, Q1 (Factorial/Div3)}\\
  3-bit values are 0 to 7.\\
  Factorial output is 1 for input values \(0,1,2\), so
  \[
  Factorial = \Sigma m(0,1,2)=A'B'C' + A'B'C + A'BC'.
  \]
  Div3 output is 1 for values \(0,3,6\), so
  \[
  Div3 = \Sigma m(0,3,6)=A'B'C' + A'BC + ABC'.
  \]

  \item \textbf{2023\_Final\_Exam\_SuppDA.pdf, Q1 (Even/Odd)}\\
  For 3-bit number with LSB \(C\): even iff \(C=0\), odd iff \(C=1\).\\
  Therefore
  \[
  E=C',\qquad O=C.
  \]
  One key difference: a multiplexer selects one input to output; a decoder maps an input code to one-hot outputs.

  \item \textbf{exam2025-main.pdf, Section 2 Q1 (Priority IRQ)}\\
  Same result as B7:
  \[
  OUT0=IRQ0,\quad OUT1=IRQ0'IRQ1,\quad OUT2=IRQ0'IRQ1'IRQ2.
  \]
  The don't-care view marks lower-priority IRQs as irrelevant whenever a higher-priority IRQ is 1.

  \item \textbf{exam24r.pdf, Q2b and Q3a (MUX and timing)}\\
  For mux-based logic questions, evaluate by select-line decoding first, then read selected data input.\\
  For timing questions, enumerate all input-to-output paths and sum gate delays; the largest sum is the critical path.
\end{enumerate}

\end{document}
